\label{app:pure-families}
Let \(T_r\) be the planar term in diagram row \(r\), including its external multiplier.
All 32 terms admit one common representation
\[
 \boxed{\displaystyle T_r=\sum_{i=1}^{48}\mathfrak{l}_{ri}(z,\bar z)\,P_i(z,\bar z).}
\]
Here \(\mathfrak{l}_{ri}\) is the complete leading-singularity prefactor in
this basis. The functions \(P_1,\ldots,P_{43}\) are a rationally linearly
independent basis of the weight-eight svMPLs occurring in the fixed external
ordering. The remaining five labels are normalized onefold primitives,
\[
 (P_{44},P_{45},P_{46},P_{47},P_{48})=(P_y,CP_y,Q,SQ,WQ),
 \qquad Q=(z-\bar z)v I_{26}.
\]
The maps \(C,S,W\) are specified below. The weight-seven integrand of \(Q\)
is denoted by \(\mathcal K_7\); its name in the ancillary Maple file is
\texttt{I26P7SV}. For example, \(P_{16}=P_v\) occurs in
\[
 \frac{I_{214}}v=\frac{P_{14}}d+
 \frac{P_{15}}{d(u-1)}+\frac{P_{16}}{d(v-1)}+\frac{uP_{44}}{d y},
 \qquad I_{37}=\frac{vP_{16}}{d(v-1)},\qquad d=z-\bar z .
\]
The two functions multiplying \(v/[d(v-1)]\) in \(I_{246}\) and \(I_{167}\)
are both \(2P_{16}\).
\begin{center}
\begin{tabular}{@{}lrrr@{}}\toprule
Comparison & svMPLs & Onefolds & Total\\\midrule
Occurrences in the 32 rows & 62 & 6 & 68\\
Equal in the displayed normalization & 52 & 5 & \textbf{57}\\
Up to a nonzero rational constant & 47 & 5 & \textbf{52}\\
Up to that constant and external crossing & 33 & 2 & \textbf{35}\\
\bottomrule\end{tabular}
\end{center}
The second and third lines keep \((z,\bar z)\) fixed. The last line identifies
functions related by one of the six permutations of the cross ratios. A rational
constant can be moved into the leading singularity. A change of basis mixing two
prefactors can also mix their two functions, so the counts refer to the stated
decomposition. Each product of lower-loop functions counts as the full function
multiplying its prefactor. The two crossed scalar onefolds in row~14 are kept as
separate evaluations.

At fixed ordering there are 47 svMPLs up to scale. Four further rational
linear relations reduce their rank to 43. For instance the pure function in
row~8 is \(P_{24}+P_{27}/3\), the sum of the square fishnet and a third of the
two-loop ladder square. Exact elimination gives a \(43\times43\) word minor
with determinant 514{,}223 modulo \(1{,}000{,}003\). Substitution of the
43 functions reconstructs all 62 svMPL occurrences over the rationals.
Crossing gives the 35 families in the last line of the table; their maps
can change the external arguments, while the 48 labels in the boxed formula
all use the same external ordering.
To compare two svMPLs we first resolve their rational prefactors exactly in
\(z,\bar z\). At each word and period monomial we solve for its coefficient on
the specified prefactors and substitute back as a rational identity.
We then transform both endpoints and all moving letters to the same chart.
The shuffle relations and the period algebra in
\href{https://github.com/oliverschnetz/HyperlogProcedures}{Schnetz's
\textsc{HyperlogProcedures} 0.8} put each function in the same word basis.
Two coefficient vectors agree up to scale precisely when their ratio at one
nonzero word agrees at every word, including pure periods.
\clearpage

\paragraph{The 32 planar integrals in the common basis.}
The middle column includes the correlator multiplier. Its prefactor in front of
\(P_i\) in the right column is the matrix entry \(\mathfrak{l}_{ri}\).
Here \(u=z\bar z\), \(v=(1-z)(1-\bar z)\), \(d=z-\bar z\) and
\(y^2=(1-u)^2-4v\), \(y_C^2=(1-u)^2-4uv\). Both roots follow the sheets of
their primitives.
\begin{center}\normalsize\renewcommand{\arraystretch}{2.05}
\begin{tabular}{@{}r@{\hspace{4mm}}l@{\hspace{5mm}}l@{}}\toprule
Row & Planar term \(T_r\) & Expansion \(\sum_i\mathfrak{l}_{ri}P_i\)\\\midrule
1 & \(I_{175}\) & \(\displaystyle \frac{1}{d}P_{1}\)\\
2 & \(I_{173}\) & \(\displaystyle \frac{1}{d}P_{2} +\frac{u - v - 1}{d^{2}}P_{3}\)\\
3 & \(I_{208}\) & \(\displaystyle \frac{1}{d}P_{4}\)\\
4 & \(I_{181}\) & \(\displaystyle \frac{1}{d}P_{5} +\frac{1}{d (u - 1)}P_{6}\)\\
5 & \(I_{174}\) & \(\displaystyle \frac{1}{d}P_{4}\)\\
6 & \(I_{176}\) & \(\displaystyle \frac{1}{d}P_{7} +\frac{1}{d (u - 1)}P_{8}\)\\
7 & \(\tfrac{1}{u}I_{233}\) & \(\displaystyle -\frac{1}{d (u - 1)}P_{9} -\frac{1}{d (v - 1)}P_{10} -\frac{1}{d (u + v - 1)}P_{11}\)\\
 & & \(\displaystyle\quad +\frac{1}{d^{2}}P_{12}\)\\
8 & \(vI_{140}\) & \(\displaystyle \frac{v}{d^{2}}P_{24} +\frac{v}{3 d^{2}}P_{27}\)\\
9 & \(I_{187}\) & \(\displaystyle -\frac{1}{d (v - 1)}P_{10} +\frac{1}{d}P_{13}\)\\
10 & \(\tfrac{1}{v}I_{214}\) & \(\displaystyle \frac{1}{d}P_{14} +\frac{1}{d (u - 1)}P_{15} +\frac{1}{d (v - 1)}P_{16} +\frac{u}{d y}P_{44}\)\\
11 & \(I_{246}\) & \(\displaystyle \frac{2 v}{d (v - 1)}P_{16} +\frac{v}{d^{2}}P_{17}\)\\
12 & \(I_{212}\) & \(\displaystyle \frac{u - v - 1}{d^{2}}P_{3} +\frac{1}{d}P_{18} +\frac{v + 1}{d (v - 1)}P_{19}\)\\
 & & \(\displaystyle\quad -\frac{u - v - 1}{2 d^{2}}P_{24}\)\\
13 & \(uI_{144}\) & \(\displaystyle -\frac{u}{d (u - 1)}P_{20} -\frac{u}{d (u - v - 1)}P_{21} +\frac{u}{d^{2}}P_{22}\)\\
14 & \(\tfrac{u}{v}I_{260}\) & \(\displaystyle \frac{u}{d^{2}}P_{23} +\frac{v - 1}{d^{2}}P_{24} +\frac{1}{d (u - 1)}P_{25}\)\\
 & & \(\displaystyle\quad +\frac{u}{d (u - 1)}P_{26} +\frac{1}{d^{2}}P_{43} -\frac{2 u}{d y}P_{44}\)\\
 & & \(\displaystyle\quad +\frac{2}{d y_C}P_{45} -\frac{2}{d}P_{47} -\frac{2}{d}P_{48}\)\\
15 & \(\tfrac{v^{2}}{u}I_{49}\) & \(\displaystyle \frac{v}{d^{2}}P_{27}\)\\
16 & \(I_{3}\) & \(\displaystyle \frac{1}{d}P_{28}\)\\
\bottomrule\end{tabular}
\end{center}
\clearpage
\paragraph{The 32 planar integrals in the common basis (continued).}
\begin{center}\normalsize\renewcommand{\arraystretch}{2.05}
\begin{tabular}{@{}r@{\hspace{4mm}}l@{\hspace{5mm}}l@{}}\toprule
Row & Planar term \(T_r\) & Expansion \(\sum_i\mathfrak{l}_{ri}P_i\)\\\midrule
17 & \(I_{5}\) & \(\displaystyle \frac{1}{d}P_{29}\)\\
18 & \(I_{39}\) & \(\displaystyle \frac{v}{d^{2}}P_{30}\)\\
19 & \(vI_{38}\) & \(\displaystyle \frac{v}{d^{2}}P_{24}\)\\
20 & \(I_{21}\) & \(\displaystyle \frac{1}{d}P_{31}\)\\
21 & \(I_{37}\) & \(\displaystyle \frac{v}{d (v - 1)}P_{16}\)\\
22 & \(\tfrac{v}{u}I_{27}\) & \(\displaystyle \frac{v}{d (v - 1)}P_{32}\)\\
23 & \(\tfrac{v}{u}I_{49}\) & \(\displaystyle \frac{1}{d^{2}}P_{27}\)\\
24 & \(I_{160}\) & \(\displaystyle \frac{1}{d}P_{1}\)\\
25 & \(I_{167}\) & \(\displaystyle \frac{2 v}{d (v - 1)}P_{16} +\frac{v}{d^{2}}P_{17}\)\\
26 & \(I_{162}\) & \(\displaystyle \frac{1}{d (v - 1)}P_{32} +\frac{1}{d}P_{33}\)\\
27 & \(I_{163}\) & \(\displaystyle -\frac{2}{d^{2}}P_{3} +\frac{2}{d (v - 1)}P_{19} +\frac{v}{d^{2}}P_{24}\)\\
 & & \(\displaystyle\quad +\frac{1}{d}P_{34} +\frac{1}{d (u - 1)}P_{35} +\frac{u}{d (u - v)}P_{36}\)\\
 & & \(\displaystyle\quad +\frac{u}{d^{2}}P_{37}\)\\
28 & \(I_{164}\) & \(\displaystyle \frac{1}{d}P_{4}\)\\
29 & \(\tfrac{u}{v}I_{133}\) & \(\displaystyle \frac{u}{d^{2}}P_{37} +\frac{u}{3 d^{2}}P_{43}\)\\
30 & \(I_{168}\) & \(\displaystyle \frac{v}{d^{2}}P_{38} +\frac{1}{d}P_{39} +\frac{v}{d (u - v)}P_{40}\)\\
 & & \(\displaystyle\quad -\frac{v}{d (v - 1)}P_{41} +\frac{2}{d}P_{46}\)\\
31 & \(uI_{169}\) & \(\displaystyle \frac{8 v}{d (v - 1)}P_{19} +\frac{v}{d^{2}}P_{42}\)\\
32 & \(\tfrac{u}{v}I_{126}\) & \(\displaystyle \frac{u}{d^{2}}P_{43}\)\\
\bottomrule\end{tabular}
\end{center}
Rows~15 and~23 contain the same two-loop ladder square \(P_{27}\) with
different prefactors. Rows~11 and~25 agree term by term. In row~27,
\(P_{19}\) appears on a simple pole and \(P_{24}\) on a double pole; their
other occurrences can be read directly from the same column. The attached
\texttt{fourloopPlanarPure.mpl} gives each \(P_i\) and the sparse list of
\(\mathfrak{l}_{ri}\). The \href{https://github.com/windfolgen/svbwalkthrough/tree/main/anc}{ancillary index}
includes its onefold dependencies and the Wolfram Language version.

\clearpage
\paragraph{All repeated functions.}
Write \(WH=H(1-z,1-\bar z)\), \(CH=H(1/z,1/\bar z)\) and
\(SH=H(z/(z-1),\bar z/(\bar z-1))\). The other maps have first arguments
\(1/(1-z)\) and \((z-1)/z\). Let \(F_{r,j}\) denote the coefficient of the \(j\)-th prefactor in Table~4 order, before reduction to the common \(P_i\) basis. The second index in row~14 follows the nine terms on the next page; the elliptic function in row~30 is fifth. The table gives one equation for every further
occurrence. Within a row of this table, \(H\) denotes the representative in the
left column. An unadorned \(H\) therefore indicates literal equality in the
same chart and normalization. Relations involving several sectors of one
integral are included.
\begin{center}\small\renewcommand{\arraystretch}{1.28}
\begin{tabular}{@{}l@{\hspace{5mm}}l@{}}\toprule
Representative \(H\) & Other occurrences in the 32 planar rows\\\midrule
\(P_{1}=F_{1,1}\) & \(F_{3,1}=F_{5,1}=F_{28,1}=-WH,\quad F_{24,1}=H\)\\
\(P_{3}=F_{2,2}\) & \(F_{14,3}=F_{19,1}=F_{27,7}=-2SH,\quad F_{27,5}=-2H,\quad F_{27,6}=-2CH\)\\
\(P_{10}=F_{7,2}\) & \(F_{9,2}=-H,\quad F_{13,1}=-WH,\quad F_{22,1}=F_{26,2}=-SH\)\\
\(P_{11}=F_{7,3}\) & \(F_{13,2}=-CH\)\\
\(P_{24}+\tfrac13 P_{27}=F_{8,1}\) & \(F_{29,1}=WH\)\\
\(P_{16}=F_{10,3}\) & \(F_{11,1}=F_{25,1}=2H,\quad F_{21,1}=H\)\\
\(P_{44}=F_{10,4}\) & \(F_{14,6}=H,\quad F_{14,7}=CH\)\\
\(P_{17}=F_{11,2}\) & \(F_{25,2}=H\)\\
\(P_{19}=F_{12,2}\) & \(F_{27,2}=-2WH,\quad F_{27,3}=2CH,\quad F_{27,4}=2H,\quad F_{31,1}=8H\)\\
\(-P_{24}+P_{43}=F_{14,1}\) & \(F_{14,2}=CH\)\\
\(P_{25}=F_{14,4}\) & \(F_{14,5}=CH\)\\
\(P_{27}=F_{15,1}\) & \(F_{23,1}=H,\quad F_{32,1}=WH\)\\
\(P_{28}=F_{16,1}\) & \(F_{17,1}=-CH\)\\
\(P_{40}=F_{30,3}\) & \(F_{30,4}=-CH\)\\
\(P_{46}=F_{30,5}\) & \(F_{14,8}=SH,\quad F_{14,9}=WH\)\\
\bottomrule\end{tabular}
\end{center}
For example, \(I_{214}, I_{37}, I_{246}\) and \(I_{167}\) all use
\(P_v\). The odd function in the last two is \(2P_v\). Their even functions
also agree with each other, and the two prefactors coincide, giving the
full magic identity \(I_{246}=I_{167}\). A different shared function links
\(I_{212}, I_{163}\) and \(I_{169}\) through the identities
\[
 F_{12,2}=\tfrac12 F_{27,4}=\tfrac18 F_{31,1}.
\]
The \href{https://github.com/windfolgen/svbwalkthrough/tree/main/runs/planarPureFunctions}{coefficient
vectors and comparison} give every match, all equivalent permutations, and
the sectors that occur only once.
\clearpage

\paragraph{Which functions survive in \(I_{260}\).}
In the four-cut benchmark let \(F\) be the odd function fitted in the reduced
graph \(M\) and \(P_0,P_u,P_v,P_y\) the four functions in \(I_{214}\), so that
\[
 \frac{I_{214}}{v}=\frac1d
 \left(P_0+\frac{P_u}{u-1}+\frac{P_v}{v-1}+\frac{uP_y}{y}\right),
 \qquad y^2=(1-u)^2-4v .
\]
Set \(D=(d\Phi_2)^2\), where \(\Phi_2\) is the native two-loop ladder, and
\(H=d^2 I_{38}\) for the square fishnet. In particular \(D=F_{15,1}=F_{23,1}\)
and \(WD=F_{32,1}\). Put \(y_C=u\,Cy\), so
\(y_C^2=(1-u)^2-4uv\), and take the primitive and its square root on the same sheet.
Define
\[
 A=F-2P_u+2(P_0-CP_0).
\]
The defining subtraction for \(R_{260}\) restores two full \(I_{168}\)
integrals in the planar term,
\[
 \frac{u}{v}I_{260}(u,v)=\frac{R_{260}(u,v)}{v}
 +\underbrace{\frac{1}{v}I_{168}(u/v,1/v)}_{x_1\leftrightarrow x_2}
 +\underbrace{I_{168}(v,u)}_{x_1\leftrightarrow x_3}.
\]
The first external permutation acts by \(S\), the second by \(W\).
Substituting \(I_{168}=R_{168}+2Q/d\) cancels both crossed remainders in
\(R_{260}\). Since \(Sd=-d/v\) and \(Wd=-d\), the surviving scalar terms are
\(-2SQ/d\) and \(-2WQ/d\), respectively. The two \(P_v\) terms of \(M\) cancel those of
\(2(I_{214}+CI_{214})\) because the fitted function is \(G=2P_v\).
In the planar normalization the result is
\begin{align*}
 \frac{u}{v} I_{260}
 &=\frac{WD-H+u(SD-H)+vH}{d^2}
   +\frac{A+u\,CA}{d(u-1)}\\
 &\quad-\frac{2u}{d y}P_y+\frac{2}{d y_C}CP_y
     -\frac{2}{d}\big(SQ+WQ\big).
\end{align*}
Reading each summand in order gives
\[
 (F_{14,1},\ldots,F_{14,9})=
 (WD-H,\ SD-H,\ H,\ A,\ CA,\ P_y,\ CP_y,\ SQ,\ WQ).
\]
This explains why row~14 shares a square-fishnet function with rows~19
and~27 and a scalar elliptic primitive with row~30. The factor of two on
the common \(I_{214}\) cut is also visible in \(M\), although \(P_v\) cancels
from the final expression. These identifications of row~14 have the same
scope as the four-cut bootstrap of \(M\).

The onefold column counts the completed primitives. For \(P_y\) the two
sheets are exchanged by the sign of \(y\); \(CP_y\) is odd in the independent
root \(y_C\). Thus the two displayed algebraic functions are distinct in the
fixed ordering. For the scalar primitive, \(CQ=-Q\) follows from inversion
of \(I_{26}\). To distinguish its three images, divide by \(d\) and take the
native OPE limit. The native \(vI_{26}\) has quartic-log coefficient \(1/24\)
at leading degree, whereas \((WQ)/d=-uI_{26}(v,u)\) starts at order
\(u\log u\). The three inversion cosets carry the natural permutation action on three points,
so equality up to scale of any two would identify all three. These
boundary values exclude that possibility. The remaining possible linear relation among the three images would be proportional to \(Q-WQ-SQ\), the alternating projection of the three images. In the upper half-plane \(Q\) has positive imaginary part and \(WQ,SQ\) have negative imaginary part by positivity of the Euclidean scalar graph. Thus the alternating projection also survives. Under crossing the five distinct
onefolds reduce to \(P_y\) and \(Q\). The weight-seven functions used inside
their last integrations occur at a different stage of the bootstrap.
